\section{Introduction}

Large language model (LLM) applications are increasingly organized as collections of specialized agents: planners, researchers, builders, critics, evaluators, and tool operators. This decomposition makes complex workflows more modular, but it also creates a routing problem. A user does not merely need ``an agent''; the user needs the right coalition of agents for the mission, with enough trust, traceability, and memory to understand why the system made a decision and whether the decision improved over time.

Many practical agent systems still rely on static routing, manually chosen roles, sequential prompt chains, or hand-authored conditional logic. These designs can be useful, but they tend to hide performance evidence. If an agent produces a hallucinated market-size claim, there is often no durable record of the failure, no transparent explanation of why that agent was selected, and no mechanism that prevents the same failure from recurring in the next run. The absence of such a trust layer becomes more severe as agent catalogs grow and as multiple agents compete for overlapping tasks.

\systemname{} addresses this gap by framing orchestration as a market-based performance exchange. Agents bid on tasks, a \marketmaker{} evaluates them with quantitative signals, selected agents execute the mission, an evaluator scores the outputs, tracing records the workflow, and persistent reputation memory updates the next routing decision. The core design is inspired by market-based task allocation in multi-agent and multi-robot systems \citep{smith1980contractnet,gerkey2004taxonomy,dias2006marketbased}, bandit exploration \citep{auer2002finite}, Bayesian reputation \citep{josang2002beta}, pairwise rating systems \citep{elo1978rating,bradley1952rank}, portfolio selection \citep{markowitz1952portfolio}, graph centrality \citep{page1999pagerank}, and modern tool-using LLM agents \citep{yao2023react,schick2023toolformer,wu2023autogen}.

The contribution of this project is not a claim that a hackathon prototype solves all agent-market design questions. Rather, \systemname{} demonstrates an end-to-end technical pattern: agent selection can be made inspectable, evaluation can become state, and orchestration can improve by routing future work through a memory-backed performance market. The live system is implemented as a Next.js application with API routes, TypeScript agent models, Gemini execution, \weave{} tracing hooks, Redis or in-memory reputation storage, mathematical routing modules, and a dashboard that exposes bids, decisions, traces, scores, and reputation changes.
